\section{Evaluaci´on Avanzada}

Se realiza una comparativa profunda de ambos modelos tras la fase de optimizaci´on, considerando tanto el rendimiento cuantitativo como la calidad cualitativa.

\subsection{Rendimiento Cuantitativo (ROUGE/BLEU)}

Se han evaluado las respuestas sobre el conjunto de test de SGD utilizando m´etricas de fidelidad de texto:

\begin{table}[H]
    \centering
    \begin{tabularx}{\textwidth}{|l|X|X|}
        \hline
        \textbf{M´etrica}          & \textbf{Qwen 3.5 (2B)}  & \textbf{Phi-4-mini (3.8B)}    \\ \hline
        ROUGE-L                  & 0.42                    & 0.39                          \\ \hline
        BLEU                     & 0.28                    & 0.25                          \\ \hline
    \end{tabularx}
    \caption{Comparativa cuantitativa sobre el split de test.}
    \label{tab:rouge_bleu}
\end{table}

\subsection{Evaluaci´on mediante LLM como Juez}

Para medir la coherencia y creatividad, se ha utilizado un modelo de mayor escala (ej. GPT-4o o Claude 3.5 Sonnet) como juez ciego para comparar las respuestas:
\begin{itemize}
    \item \textbf{Fidelidad:} Qwen 3.5 mantiene una mayor fidelidad a los nombres de servicios y APIs.
    \item \textbf{Naturalidad:} Phi-4-mini es preferido por su tono conversacional m´as natural y menos rob´otico.
\end{itemize}

\subsection{Interpretabilidad}

Se ha analizado el peso de los tokens mediante mapas de atenci´on, observando que ambos modelos centran correctamente la activaci´on en los nombres de servicios (\textit{restaurants}, \textit{hotels}) y en los par´ametros de reserva (\textit{time}, \textit{location}), validando que la arquitectura Transformer est´a procesando las entidades clave del diálogo.
